\chapter{Introduction}
\label{chap:introduction}

\section{Background to the Study}
Public transportation remains a critical component of mobility in Uganda, particularly in rapidly growing urban centers such as Kampala \cite{mak2026kampala}. While traditional taxi systems (matatus) serve a large portion of the population, there has been a growing demand for more reliable, comfortable, and organized transport alternatives \cite{rugamba2020lowcarbon}. Private shuttle services have emerged to fill this gap by offering scheduled trips, reduced overcrowding, and improved travel experiences. Despite these advantages, many of these services still operate with limited technological support \cite{bmau2025transport}.

Private shuttle operators often rely on manual processes such as phone calls and informal record-keeping to manage bookings, communicate with passengers, and coordinate trips. This approach presents several challenges, including double bookings, customer no-shows, difficulty in tracking seat availability, and inefficiencies in managing routes and schedules \cite{masinde2024uganda}. Additionally, customers may face challenges in accessing these services, as they are often only known through referrals or local advertisements, limiting their reach and accessibility.

The lack of a centralized digital platform also affects courier services offered by these operators. Package delivery coordination, pricing, and tracking are typically handled manually, making it difficult to provide real-time updates and efficient service \cite{itd2026couriercomparison}. This results in frequent customer inquiries, operational delays, and missed opportunities for optimizing vehicle usage.

With the increasing penetration of smartphones and mobile internet in Uganda \cite{ucc2024sectorreport}, there is a significant opportunity to enhance the operations of private transport services through technology. Mobile applications that support features such as seat booking, digital payments, real-time tracking, and automated communication can greatly improve both customer experience and operational efficiency \cite{vocal2026busbooking}, \cite{sayyed2024mobilecx}.

This project proposes the development of a mobile application, Traxis, aimed at supporting private shuttle and courier service providers. The system will enable users to easily discover available services, book seats, select preferred seating positions, and track journeys or deliveries in real time. For operators, the application will provide tools for managing bookings, monitoring vehicle movement, tracking revenue, and optimizing capacity utilization. By introducing a structured and user-friendly digital platform, the system aims to improve reliability, accessibility, and overall service delivery within the private transport sector in Uganda.

\section{Problem Statement}
Private shuttle transport services in Uganda, particularly in urban centers like Kampala, face significant challenges related to coordination and service management despite offering a more reliable alternative to traditional public transport. These challenges result in inefficiencies such as limited service visibility, difficulties in accessing booking information, customer no-shows, and underutilization of available seating capacity. The reliance on manual processes, such as phone calls and informal scheduling, further complicates operations and limits the ability of service providers to effectively manage passenger demand.

The absence of a structured, technology-enabled platform to support real-time booking, seat allocation, and service coordination exacerbates these issues. Consequently, customers often experience inconvenience and uncertainty when trying to access these services, while operators face operational inefficiencies that can lead to revenue losses and reduced service quality. Without an accessible and user-friendly digital solution, these challenges will persist, highlighting the need for an application such as Traxis to enhance communication, improve transparency, and optimize the management of private shuttle transport services in Uganda.
\vspace{1cm}
\section{Objectives}
\subsection{Main Objective}
% The main objective of this project is to develop a mobile application that reduces waiting times and increases the average taxi occupancy rates through real-time tracking and scheduled pickups.

The main objective of this project is to design and develop a mobile application (Traxis) that facilitates efficient management of private shuttle and courier transport services. The system will incorporate features such as real-time vehicle tracking, digital seat booking, and passenger scheduling for pickups. This initiative aims to address key operational inefficiencies by improving coordination between service providers and customers, reducing booking-related challenges such as no-shows, and optimizing seat utilization. Ultimately, the application seeks to enhance the reliability, accessibility, and overall efficiency of private transport services.
% 3
% \newpage
\subsection{Specific Objectives}
We aim to achieve the following specific objectives with the project:
\begin{itemize}
    \item[i.]  To gather comprehensive data on passenger requirements, operational challenges faced by private shuttle and courier service providers, and the functionalities of existing transport management systems.
    % To study the existing or related systems in order to ascertain the need for the proposed mobile application.
    
    \item[ii.]To synthesize and analyze the acquired data to define a comprehensive set of system requirements for the Traxis application.
    % To analyse data from taxi personnels as well as the desired customers for a comprehensive set of system requirements for the Traxis application.
    
    
    \item[iii.] To design and implement the proposed mobile application with all core feautures.
    
    \item[iv.] 
    % To test and validate the system's performance and usability.
    To test and validate the performance and overall usability of the deployed system against the defined requirements and user expectations.
    
\end{itemize}
\vspace{1cm}
\section{Scope of the Project}
The scope of the Traxis project encompassed the development of a mobile application designed to improve the organization and reliability of private shuttle and courier transport services in Uganda, with initial focus on urban-to-upcountry routes such as Kampala to Kabale. The application provides core functionalities including real-time vehicle tracking, seat booking and allocation, and passenger scheduling for pickups. It also supports basic courier service management, including request placement and delivery tracking.

The system excluded other forms of transport such as boda bodas and standard public taxi (matatu) services in its initial phase in order to maintain a clear focus and ensure manageable system development within the project timeframe.

The system caters primarily to private shuttle operators and passengers who own smartphones and had access to mobile data, with an emphasis on ease of use to accommodate a range of digital literacy levels. The application enables shuttle operators to manage bookings and identify passenger demand hotspots, while allowing passengers to make informed travel decisions based on seat availability and scheduled departures.

Additionally, the project scope covered backend development to support real-time data processing, basic communication features for operator-passenger interactions, and user feedback mechanisms to allow iterative improvement. The project did not initially incorporate advanced features such as extensive multi-city coverage; however, it was designed to allow potential future scaling based on user needs and technical feasibility. 

\section{Significance of the Study }
The Traxis project holds a substantial significance for multiple stakeholders within Uganda’s private shuttle transport ecosystem and represents a meaningful step toward digital transformation in the informal transport sector. This study was justified by several critical factors:

\begin{itemize}
    \item \textbf{To Passengers:} The application addresses customer challenges experienced in private shuttle transport services by enabling structured seat booking, real-time vehicle tracking, and scheduled pickups. It reduces uncertainty in travel arrangements by allowing users such as passengers to view available seats, select preferred seating, and receive instant booking confirmations. The system also improves coordination during journeys by updating seat availability in real time when passengers disembark along the route, allowing better utilization of vehicle capacity.

    \item \textbf{For courier services:} the application supports convenient package pickup requests, real-time tracking of deliveries, and automated status notifications to both senders and recipients. This reduces the need for manual follow-ups and improves transparency in goods transportation.
    
    \item \textbf{To private shuttle operators:} The application offers private shuttle operators improved operational efficiency through enhanced visibility of passenger demand, better booking coordination, and reduced empty return trips. It enables optimized seat utilization, minimizes losses from unfilled return journeys, and improves daily earnings. The system also helps manage challenges such as customer no-shows and manual booking overloads through structured digital confirmations.
    In addition, the application provides real-time vehicle monitoring, allowing operators to track whether a shuttle is moving, stopped, or delayed along its route. This improves operational control, enhances accountability for drivers, and enables better communication with customers regarding trip progress and estimated arrival times.
    
    \item \textbf{To Transportation System Development:} The project contributes to the gradual digital transformation and improved organization of Uganda’s private shuttle transport sector. By introducing structured booking, tracking, and coordination mechanisms, Traxis supports the development of a more reliable and efficient transport ecosystem.
    
    \item \textbf{To Technological Adoption:} This study demonstrates the feasibility of implementing mobile-based transport solutions in environments with varying levels of digital literacy. By incorporating phone-number-based access, simple interfaces, and optional customer support via calls, the system reflects appropriate technology design suited to the Ugandan context.
    
    \item \textbf{To Future Research and Development:} The project establishes a foundation for further innovation in private transport and logistics systems, including expansion to multiple routes, improved revenue analytics, integrated communication systems, and enhanced logistics tracking for both passengers and goods.
\end{itemize}
% \newpage
\section{Justification}
The justification for undertaking the Traxis project stems from the urgent need to address coordination and operational inefficiencies within Uganda’s private shuttle and courier transport services. The current largely manual system, which relies on phone calls, informal scheduling, and limited digital support, results in challenges such as booking inconsistencies, customer no-shows, underutilized vehicle capacity, and lost revenue opportunities for service providers. These inefficiencies reduce overall service reliability and limit the growth potential of the sector.

At the same time, the increasing penetration of mobile technology and internet access in Uganda presents a timely opportunity to implement a digital solution tailored to this context. The Traxis application is therefore justified as a practical and context-aware intervention designed to bridge these operational gaps. By leveraging widely available technologies such as mobile data, GPS tracking, and mobile money integration, the project aims to improve booking coordination, enhance transparency in service delivery, and support real-time communication between operators and customers.

Ultimately, the system is intended to deliver immediate and practical benefits to passengers and private shuttle operators while also improving the efficiency of courier service operations and contributing to the gradual digital transformation of the transport sector in Uganda.